\odiseachapter{matanza}{La matanza de los pretendientes}\label{cap:matanza}

Al contrario que en la \emph{Nodisea}, el original no hace un melodrama de la muerte de Antínoo, ni falta que hace. De hecho, es el primero que cae:

\begin{versocita}
Y enderechó, esto diciendo, a Antínoo una amarga saeta.\\
Estaba aquel levantando en eso una copa muy bella\\
de oro y ansas gemelas ---la asían sus manos apenas---\\
para su vino beber, y a matanza su alma era ajena.\\
¿Cómo iba allí uno a pensar que entre aquellos sentado a la mesa\\
un hombre, solo entre muchos ---por más que sobrado de fuerza---,\\
la mala muerte vendría a traerle y el ángela negra?
\end{versocita}

Odiseo lo despacha mientras Antínoo da un trago de vino, sin darle tiempo a rechistar. Como es marca de la casa, la descripción de su muerte es exacta y terrible:

\begin{versocita}
A él encarando, la gorja Odiseo le abrió con la flecha,\\
y por el mórbido cuello salió la punta a derechas.\\
Venciose a un lado, y la copa cayó de sus manos sin fuerzas,\\
y un caño espeso bajaba por su nariz de carrera\\
de humana sangre, y con un golpe brusco del pie echó la mesa\\
lejos de sí, y el manjar que tenía esparció por la tierra:\\
el pan, las carnes doradas en sangre y polvo. Se encrespan\\
los pretendientes, al verlo caer, en barbulla violenta,\\
y saltan de sus sitiales, y en pánico por la tarbea\\
miran, buscando, alredor los firmes muros de esa;\\
mas que tomar no había escudo ninguno ni lanza cereña.\\
Y con palabra enconada así a Odiseo lo increpan:
\end{versocita}

\begin{versocita}
«A malas has disparado, extranjero; pero otro torneo\\
no te vendrás a encontrar: ya has ganado el barranco funesto.\\
Pues ¿qué?, si has muerto al que era el más preclaro mancebo\\
de Ítaca belmoceril: de los buitres serás alimento».
\end{versocita}

Todavía no se han dado cuenta de la que les espera. El torvo héroe los saca de su error enseguida:

\begin{versocita}
«Perros, pensabais que nunca de vuelta yo a casa llegara\\
de aquella tierra troyana, y así mi casal devorabais,\\
y por la fuerza os entrabais al lecho con mis esclavas,\\
y a mi mujer en asedio ---y yo vivo--- la cortejabais\\
sin miedo a los que en el cielo morada tienen bien ancha\\
ni a que de hombre viniera andando el tiempo venganza:\\
ahora atado lleváis el nudo de la matanza».
\end{versocita}

Eurímaco, entonces, intenta salir del apuro echándole la culpa de todo al caído Antínoo:

\begin{versocita}
«Si es la verdad, Odiseo de Ítaca, que has regresado,\\
razón te sobra en lo dicho de cuanto hicimos los dánaos,\\
desmanes mil en tu casa, y mil también en el campo.\\
Mas el que culpa ha de todos lo tienes ahí derrumbado,\\
Antínoo; que era este solo el que nos empujaba a esos casos,\\
no porque boda anhelara o necesitárala tanto,\\
sino mirando otro fin que no quiso el Cronión consumarlo:\\
ser en Ítaca él, la de firme solar, soberano\\
luego de armarle a tu hijo celada para matarlo.\\
Él yace ahora en su suerte, perdona tú a tus vasallos;\\
que del común ya te haremos arreglo después en reparo\\
de cuanto se haya bebido y comido aquí en tu palacio,\\
y cada uno un desquite en veinte bueyes tasado\\
te pagaremos en bronce y en oro hasta haber aplacado\\
tu corazón: y hasta aquí no es baldón que tú sigas airado».
\end{versocita}

Se trata de una escena que Nolan ---es decir, Hollywood--- nos escamotea. ¿Por qué? Porque lo razonable, \emph{lo honorable}, desde nuestra óptica moderna, sería perdonar a los enemigos. A fin de cuentas, Odiseo se ha vengado del traidor, ha dado un escarmiento y le ofrecen compensación por los desmanes cometidos.

Nada de eso. Al héroe nadie va a privarlo de su venganza:

\begin{versocita}
Y a él, mirándolo torvo, le dijo el ardido Odiseo:\\
«Eurímaco, ni aunque pagarais con todos los bienes paternos,\\
cuantos ahora tenéis, y hasta más que añadierais a esos,\\
ni así pondría a mis manos de la matanza yo el freno\\
hasta que la demasía pagado haya bien el cortejo.\\
Pelear de frente es ahora lo que a vosotros os dejo\\
o huir si a muerte y malhado alguno ve escapamiento;\\
mas se me da que ninguno va a hurtarse del trance funesto».
\end{versocita}

Eurímaco se da cuenta de que lo único que les queda es pelear o morir y dice:

\begin{versocita}
«Amigos, no detendrá este hombre su mano en holganza,\\
que ahora que dueño se ha hecho del arco pulido y la aljaba\\
desde el umbral reluciente nos flechará hasta que haga\\
carnicería de todos: pensemos ya en la batalla.\\
Desenvainad las espadas, poned vuestras mesas de cara\\
a las saetas de pronto morir, y vayamos en masa\\
todos a él, y de umbral y de puertas hagamos que salga\\
para llegarnos por la ciudad y que vuele la alarma;\\
así ya ese hombre hoy por última vez con el arco tirara».
\end{versocita}

Atención, los pretendientes no necesitan que el cabrero les tire las armas desde el piso de arriba, como ocurre en la película. Ya están armados, llevan a cuestas sus espadas de bronce y, en ese sentido, Homero es más atrevido ---por enésima vez--- que Nolan. El tímido director pretende que los dánaos no tienen apenas con qué defenderse, lo que hace la hazaña de Odiseo menos espectacular ---aunque quizás más creíble a ojos de un espectador moderno---. El Poeta, sin embargo, no solo los arma, sino que les da la voluntad de combatir. Y, la verdad, uno se pregunta cómo se las va a arreglar el furioso arquero para despachar a un centenar de jóvenes blandiendo espadas y protegiéndose de sus flechas con improvisados escudos de madera. La lectura, entre líneas, está muy clara. La hazaña solo es posible porque los pretendientes son humanos y Ulises, aquí, tiene algo ---mucho--- de divino. Y no es un dios misericordioso. Odiseo, el de las muchas caras, es también el dios de la muerte y la destrucción.

Eurímaco, de hecho, muere con honor, intentando defenderse:

\begin{versocita}
A la que dijo él así, sacó su espada buída\\
a filo y filo aguzada, broncina, y echósele encima\\
a hórrido grito a la vez que el divino Odiseo una vira\\
le disparaba alcanzándole el pecho bajo la tetilla:\\
raudo en el hígado el tiro se hundió; y la espada que asía\\
se le escurrió entre los dedos, y desmadejado él caía\\
enguruñido sobre una trabanca, y rodó la comida\\
y el cáliz doble de cuenco; y en estertor de agonía\\
golpeaba el polvo su frente, y por detrás sacudía\\
con ambos pies el sitial; y la niebla veló sus pupilas.
\end{versocita}

El siguiente en abrazar la parca es Anfínomo, el sensato pretendiente a quien el mendigo Odiseo había aconsejado que se marchara a casa, y no es este quien lo mata, sino Telémaco:

\begin{versocita}
Lanzose entonces Anfínomo hacia el glorioso Odiseo\\
bien a derecho ---en su mano llevaba el alfanje malfiero---\\
por si apartarlo de puertas pudiera. Mas antes Telémaco\\
fue por detrás y, arrojándolo, hincole el venablo bronceño\\
entre hombro y hombro a mitad, y se lo pasó por el pecho:\\
él al caer retumbó, y con toda su frente dio al suelo.
\end{versocita}

Inmediatamente, el hijo de Odiseo corre junto a su padre y le dice:

\begin{versocita}
«Padre, un escudo y dos lanzas ahora voy a traerte\\
y yelmo todo de bronce bien ajustado a tus sienes,\\
y me armaré yo también, y a porquero y boyero de bueyes\\
daré sus armas: ceñirnos los cuatro ahora bien nos conviene».
\end{versocita}

Así que Homero no pretende que sea Ulises quien luche en solitario, pues le acompañan en la refriega su hijo y sus fieles. El héroe, que sabe lo suyo de batallas, le dice:

\begin{versocita}
«Tráelas corriendo mientras me defienden las flechas que quedan,\\
no vaya a ser que al dejarme aquí solo me ganen la puerta».
\end{versocita}

\begin{versocita}
Tal dijo; y obedeciendo el muchacho a su padre, se marcha\\
al almacén al que habían llevado las ínclitas armas.\\
De allí sacó cuatro escudos, y cuatro parejas de lanzas,\\
y cuatro yelmos de bronce con sus cimeras crinadas;\\
con ellos raudo salió, y raudo al padre llegaba.\\
Él mismo el bronce a su cuerpo ciñó muy antes que nada,\\
y, como él, los dos siervos vistieron sus armas galanas,\\
y del ardido Odiseo a redor, mienteviva, se plantan.
\end{versocita}

Hay una escena de la matanza de Nolan que da casi risa. Matt Damon se defiende sin más arma que una flechita, gesticulando mucho para apartar a los enemigos que le acosan, tan pánfilos como él. Homero nos pinta una batalla mucho más fiera:

\begin{versocita}
Mientras duráronle a este para defenderse los dardos,\\
de los galanos fue haciendo blanco en la sala tras blanco,\\
y siempre a uno mataba, y caían amontonados.\\
Mas cuando ya al rey flechero las flechas lo abandonaron,\\
del firme muro en la sala contra un pilar apoyado,\\
por la crujía de todo esplendor dejó quieto el arco,\\
luego a los hombros echose escudo cuatriadobado,\\
caló en su testa robusta el yelmo bienaforjado,\\
corcelicrino ---de arriba pendía espantable penacho---,\\
y, rematados en bronce, asió dos recios venablos.
\end{versocita}

¿Cómo te sentirías, asustada lectora, petrificado lector, si vieras delante de ti a un gigantón furioso, enfundado en bronce, armado hasta los dientes? Pero en esto, el cabrero Melantio consigue hacerse con no pocas lanzas y escudos, lo que desequilibra todavía más el ya desigual combate:

\begin{versocita}
Y el cabrerón de las cabras, Melantio, marchó, esto diciendo,\\
por las lumbreras traseras al almacén de Odiseo.\\
Sacó de allí doce escudos, y mismas lanzas con ellos\\
y mismos fueron y en bronce corcelicrinos los yelmos.\\
Y vuelve a todo correr por repartirlo al cortejo.\\
Le flaquearon las piernas y el corazón a Odiseo\\
a la que vio que, ya de armas ceñidos, estaban blandiendo\\
largo rejón en sus manos: trabajo iba a ser no pequeño.
\end{versocita}

Cuando Melantio se dirige al almacén para seguir trayendo armas, Ulises manda a Eumeo y Filetio que lo intercepten. Sin embargo, no les da orden de matarlo, sino de dejarlo atado de pies y manos y colgado del techo ---las razones ya nos las imaginamos---. Así que no es Telémaco, como en la \emph{Nodisea}, quien enfrenta al cabrero, sino los dos fieles sirvientes, que, tras cumplir su misión, se unen al héroe y su hijo. La situación, entre tanto, se ha puesto fea. Son cuatro contra cien.

\begin{versocita}
Ellos vistieron sus armas, cerraron la puerta brillante,\\
y hacia el sagaz Odiseo volvieron, el milmañasartes.\\
Estaban cuatro, alentando furor, allí en los umbrales,\\
y enfrente, dentro de casa, turba de muchos pujante.
\end{versocita}

Así las cosas, es inevitable que Atenea aparezca:

\begin{versocita}
Atena, la hija de Zeus, a aquellos vino a allegárseles\\
a Méntor siendo en figura y a par en voz semejante.\\
Y se alegró de esa vista Odiseo y le dijo al instante:
\end{versocita}

\begin{versocita}
«Méntor, ayuda a este estrago, recuerda a tu amigo querido,\\
quien tantos bienes te hacía, el par en años contigo».\\
Dijo, aun sabiendo que era Atena la Salvadora.
\end{versocita}

Uno de los pretendientes, Agelao, sin saber que tiene delante a una diosa, amenaza al supuesto Méntor ---sin darse cuenta tampoco, aparentemente, de que ha aparecido de la nada---:

\begin{versocita}
Enfrente los del cortejo gritaban en una voz sola;\\
y a ella le lanza reproche el primero Agelao Damastórida:\\
«Méntor, que no te engatuse con ese su verbo Odiseo\\
para luchar contra los pretendientes y a él socorrerlo.\\
Que nuestra empresa vendrá a terminar como ahora te cuento:\\
tan pronto como a esos dos ---al padre y al hijo--- matemos,\\
muerto con ellos caerás tú también, que es muy caro el empeño\\
que te propones aquí; tu cabeza será nuestro precio.\\
Y cuando a bronce os hayamos quitado las fuerzas y aliento,\\
cuantas habiendas posees ---las que están bajo techo y a cielo---\\
las juntaremos con las de Odiseo; y ni dejaremos\\
en tu palacio vivir a tus hijos como herederos,\\
ni que tu viuda y tus hijas por Ítaca anden sin dueño».
\end{versocita}

Atenea se enfada con estas palabras, pero no con los galanes, que considera carne de cañón, sino con su protegido Odiseo, por pusilánime. Claramente los estándares de los olímpicos no son como los de los humanos:

\begin{versocita}
«Ya en ti no están, Odiseo, el brío firme y pujanza\\
de cuando tú por la bien blanca Helena, la bienempadrada,\\
por nueve años luchaste contra los teucros sin tasa\\
y a muchos hombres mataste en la penosa batalla\\
y por tu ardid se tomó la ciudad de su rey calliclara.\\
¿Cómo es que ahora, cuando has alcanzado tu hacienda y tu casa,\\
plantado ante estos galanos, tan fuerte tú, te amilanas?\\
Ah, tente aquí, melón dulce, a mi par y mira mis ganas,\\
que vas a ver cómo son ante cáfila hostil las agallas\\
con las que paga este Méntor Alcímida gracias pasadas».
\end{versocita}

Y en otro de sus giros desconcertantes, la diosa, después de semejante bravuconada, se convierte en golondrina y se instala en una viga del techo desde donde seguir cómodamente la refriega. A ojos humanos, simplemente, Méntor desaparece. Pero los aqueos, enzarzados en mortal combate, no están para trucos de magia. Agelao ve su oportunidad de ganar la batalla:

\begin{versocita}
Buscó Agelao Damastórida en los pretendientes arrojo,\\
y Anfimedonte también, y Eurínomo, y aun Demoptólemo,\\
y el de Políctor, Pisandro, y el ardentísimo Pólibo,\\
a la sazón los más bravos, muy por encima de todos,\\
de los que, vivos aún, por sus vidas peleaban; que a otros\\
el arco y ristra de dardos rodar los hizo ya al polvo.\\
Y fue Agelao quien tomó la palabra y habló de este modo:
\end{versocita}

\begin{versocita}
«Amigos, pronto tendrá este hombre su mano en holganza;\\
Méntor ya veis que no está, que dijo sus baladronadas\\
y los dejó en los lumbrales igual de solos que estaban.\\
Así que no tiréis todos la larga lanza a la brava,\\
que tiren seis los primeros, a ver si Zeus nos regala\\
herida hacer a Odiseo y prez ganar de la fama.\\
No hay de los otros cuidado a la que a tierra ese caiga».
\end{versocita}

Agelao tiene muy claro que, si abaten a Odiseo, la batalla se acaba. Seis de los pretendientes le arrojan entonces las lanzas, pero Atenea las desvía.

\begin{versocita}
Dijo, y lanzaron entonces los seis, como había mandado,\\
en ansia viva; e inanes Atena dejó esos venablos.\\
Uno en la jamba no más de la bien firme sala hizo blanco,\\
otro en la puerta de sólido engarce, y de otro fue el dardo\\
broncicebado a golpear en el muro y caer rebotando.
\end{versocita}

Hay dos lecturas posibles aquí. Una es admitir que incluso un asesino profesional como Odiseo no podía derrotar a cien jóvenes armados y, en consecuencia, debe su victoria exclusivamente a la intervención divina. Es una interpretación que recorre, como hemos visto, toda la \emph{Odisea} y, en general, toda la épica griega. Los dioses, en último término, son los que deciden y, si Ulises sale victorioso de la refriega, no es por valiente, hábil o astuto, sino por ser el protegido de Atenea.

Otra lectura, que me gusta más, es asumir que Homero usa a menudo a los dioses en sentido metafórico. Zeus nos da una pista durante el concilio de los olímpicos nada más arrancar la \emph{Odisea}:

\begin{versocita}
Ay, los humanos siempre andan echando a los dioses la culpa;\\
que es que les van de nosotros sus males repiten, y nunca\\
ven que no estaba en su sino el revés por sus propias locuras.
\end{versocita}

A la inversa, quizás los humanos también «echan la culpa» a los olímpicos cuando la suerte les favorece. Puede que no haga falta recurrir a la diosa ojiglauca para frenar las lanzas. Tal vez, simplemente, Odiseo es un tipo afortunado, como se dirá a sí mismo, años más tarde, viejo y desengañado en Ítaca:

\begin{versopropio}
Al menos ahora eres lo bastante lúcido\\
para no engañarte, como solías hacer\\
cuando te imaginabas ser el gran Ulises,\\
el héroe y el viajero y el amante de Circe.\\
Cuando recuerdas, todo lo que ves\\
es ruido y furia. Un matón cruel y astuto\\
que tuvo más suerte que los demás.
\end{versopropio}

El caso es que las lanzas no aciertan y ahora es el turno de Ulises y los suyos:

\begin{versocita}
Cuando los cuatro las lanzas de los del cortejo esquivaron,\\
habló el sufrido Odiseo a los que tenía a su lado:\\
«Diría, amigos, que es vez de que pica arrojemos nosotros\\
sobre el tropel de los pretendientes que están tan ansiosos\\
de coronar anteriores oprobios con nuestro despojo».
\end{versocita}

Por supuesto, ni Odiseo ni sus fieles fallan.

\begin{versocita}
Dijo; arrojaron entonces sus lanzas vibrantes ya ellos\\
tirando al frente; e hincó a Demoptólemo muerte Odiseo,\\
a Euríades muerte Telémaco, a Élato muerte el porquero,\\
e hincó el boyero de bueyes muerte a Pisandro. Mordieron\\
al mismo tiempo los cuatro allí el sinfínido suelo,\\
y recularon buscando un rincón los demás del cortejo.\\
Y ellos lanzáronse ahí a arrancar el astil de los muertos.
\end{versocita}

Los pretendientes tiran de nuevo, y esta vez Telémaco y el porquero reciben heridas leves:

\begin{versocita}
Los pretendientes sus lanzas por vez segunda arrojaron\\
en mismas ansias, y Atena frustró casi todos sus dardos.\\
Uno en la jamba otra vez de la bien firme sala hizo blanco,\\
otro en la puerta de sólido engarce, y de otro el venablo\\
dio, ponderoso de bronce, en el muro y cayó rebotando.\\
Anfimedonte, no obstante, a Telémaco hiriolo en la mano,\\
que le rozó la muñeca, y le hizo por cima desgarro.\\
Ctesipo a Eumeo pasole por sobre el escudo su dardo\\
que, haciendo araño en su hombro, siguió hasta el suelo volando.
\end{versocita}

Inevitable imaginar ahora a Telémaco llegando a casa con la muñeca vendada y diciéndole a su madre: «deberías haber visto al otro tipo».

\begin{versocita}
Y ya otra vez Odiseo sagaz y su corro arrojaron\\
sus azagayas punteras sobre el tropel de galanos.\\
Asolador Odiseo en Euridamante hizo blanco,\\
Telémaco en Anfimedonte, en Pólibo el guardamarranos,\\
y por su parte a Ctesipo el que tiene de bueyes cuidado\\
le hincó la suya en el pecho mientras decíale ufano:
\end{versocita}

\begin{versocita}
«Politersida, el amigo de befas, no más te abandones\\
a altisonancias de poco saber, es mejor que a los dioses\\
dejes el verbo, que pueden los dioses muy más que los hombres.\\
De aquella pata de vaca adehala es en torna este golpe,\\
la que le diste a Odiseo mendigo en sus propios salones».
\end{versocita}

Filetio se refiere a lo que ha ocurrido un rato antes, durante la cena, cuando Ulises aparece disfrazado de mendigo. Ctesipo, hijo de Politerses (de ahí «Politersida»), le arroja una pata de vaca al supuesto pedigüeño, que este esquiva sin esfuerzo. El boyero le devuelve la broma con un lanzazo en el pecho.

A partir de aquí, las cosas se ponen feas para los pretendientes. El pánico cunde entre ellos y les hace romper filas. La refriega se convierte en carnicería:

\begin{versocita}
[...] Al punto Odiseo\\
al de Damástor herida de lanza le abrió en cuerpo a cuerpo.\\
Luego Telémaco a Leócrito hirió, al hijo de Evénor,\\
con su azagaya en mitad de la ijada y de claro entró el rejo:\\
aquel de bruces cayó, y con toda su frente dio al suelo.\\
Entonces ya alzó Atenea desde lo alto del techo\\
su adarga malamuertera, y el pánico se hizo aleteo.\\
E iban de acá para allá como vacas que, al pasto saliendo,\\
zumbón el tábano espanta metiéndose en ellas molesto,\\
por las sazones de abril, cuando largos los días van siendo.\\
\stanzabreak
Como los buitres zarpirrecorvos y piquirretuertos\\
se echan encima de las avecicas desde los oteros\\
y ellas se encogen bajo las nubes y a ras van de suelo\\
y ellos saltando sobre ellas las matan, pues ni abrigadero\\
para esas hay ni escapada, y el cobro alegra al montero,\\
así, irrumpiendo en la sala, los cuatro a los del cortejo\\
van golpeando a voleo en redor, vuelan gritos horrendos\\
de las cabezas quebradas, y humea de sangre ya el suelo.
\end{versocita}

Aquí, Homero y Tobal resuenan. La imagen de los buitres «zarpirrecorvos y piquirretuertos» cayendo sobre las avecicas nos hiela la sangre en las venas y a la vez nos deja maravillados de las posibilidades del lenguaje. Que en estos tiempos en los que WhatsApp impone vocablos cada día más pobres se nos ofrezca una palabra que condensa «corvas zarpas» y otra que dibuja de una manera tan gráfica los picos de los buitres sería suficiente justificación para leer ---y pedirles a nuestros hijos que leyeran--- esta traducción de la \emph{Odisea}.

Aunque hay que reconocer que el Poeta no es lectura apta para corazones tiernos. La violencia que condensan los versos es sobrehumana: «vuelan gritos horrendos / de las cabezas quebradas, y humea de sangre ya el suelo». Las siguientes estrofas subrayan ---por si no lo teníamos ya claro--- la crueldad del héroe:

\begin{versocita}
Leodes, corriendo hacia él, las rodillas tomó de Odiseo,\\
y suplicándole así le decía en alígeros verbos:\\
«De hinojos ruego, Odiseo, y tú tenme pena y reparo:\\
a fe que nunca a ninguna mujer de las de tu palacio\\
ofensa dije ni hice, que siempre a los otros galanos,\\
si alguno cosa así hacía, trataba yo de frenarlos.\\
Mas ni con eso logré sujetar la maldad de sus manos,\\
y así, por su altanería, funesto sino alcanzaron.\\
Yo, que un arúspice he sido entre ellos y nada hice malo,\\
¿he de caer? ¿Del bien hecho hasta ahora no hay gracia ni pago?».
\end{versocita}

Recuerda, memoriosa lectora, sesudo lector, que Leodes, el de las finas manos, es un buen tipo, pues el propio Homero nos dice que era el único a quien las iniquidades resultaban odiosas y que se indignaba con todos los pretendientes por gorrones y maleducados. ¿Lo perdonará Ulises?

\begin{versocita}
Y lo miró torvifiero el ardido Odiseo y le dijo:\\
«Ya que te jactas de haber sido tú su oficiante adivino,\\
seguro, sí, que mil veces habrás en la sala pedido\\
porque mi dulce regreso lejos quedara, y contigo\\
mi esposa amada se fuera y te alumbrara a ti hijos;\\
por eso no escaparás a este final maldolido».
\end{versocita}

\begin{versocita}
Tal que hubo dicho, una espada tomó entre sus túrgidos dedos\\
que allí yacía ---la había soltado Agelao por el suelo\\
cuando murió---, y con ella segole a mitad el pescuezo.\\
Y su cabeza rodó por el polvo, aún balbuciendo.
\end{versocita}

Nada que ver con el asaeteado Cristodiseo que nos vende Nolan. Todo personaje de la literatura puede reinterpretarse si se quiere ---la literatura quizás no sea otra cosa que escribir y reescribir siempre las mismas historias--- pero hay que saber a qué se juega. Cuando uno se empeña en vestir un lobo con piel de oveja, el riesgo es caer en la caricatura.

Terminada la matanza, el lobo en cuestión recorre la sala, por si queda alguien por rematar.

\begin{versocita}
Miró Odiseo por toda la estancia por si un hombre vivo\\
hallaba allí agazapado esquivando su muy negro sino.\\
Mas vio que ya estaban todos en sangre y polvo caídos.\\
Muchos. Que igual que esos peces que dan a cala en abrigo\\
los pescadores sacándolos fuera del mar espumizo\\
a la su red enmallados, y sobre la arena esparcidos\\
quedan, echando de menos las olas del mar, reunidos\\
mientras les quita la vida el haz del sol con su brillo,\\
así yacía el tropel de galanos, en parva abatidos.
\end{versocita}

Cuando la anciana Euriclea entra en la habitación, el Poeta nos muestra en qué se ha convertido su bebé:

\begin{versocita}
Halló a Odiseo en mitad de los cuerpos que ha nada abatieran\\
embadurnado de sangre y matanza cual si un león fuera\\
que de comer hasta hartarse de un toro campero volviera:\\
el pecho todo ---y a un lado y a otro sus dos carrilleras---\\
rojo de sangre, espanto de ver al mirar de cualquiera;\\
tal, pringoriento Odiseo por brazo arriba y por piernas.
\end{versocita}

Imagina, asustada lectora, trémulo lector, que eres tú quien contempla al criminal embadurnado de sangre. Posiblemente saldrías corriendo, caerías de hinojos entre arcadas o te desmayarías. Euriclea, por el contrario, se pone a gritar de alegría.

\begin{versocita}
Apenas vio los cadáveres ella y la sangre sin cuenta,\\
rompió a gritar de alegría, que vio muy grande proeza,\\
pero Odiseo del ansia la tuvo y su gana refrena,\\
y a ella, alzando la voz, palabra en alas le lleva:\\
«Por dentro alégrate, vieja, y el grito contén de alegría;\\
gloriarse sobre los muertos ni es santo ni va a ley divina».
\end{versocita}

Cuando Odiseo sea tan anciano como Euriclea, los pretendientes asesinados se sumarán a las hordas de fantasmas que lo acosan:

\begin{versopropio}
Pero lo más vívido de todo es ver los rostros\\
de esos locos muchachos, esa alegre pandilla de borrachos,\\
que tan solo querían divertirse a tu costa.\\
Es cierto, se comieron tus corderos,\\
se bebieron tu vino, pretendiendo cortejar a tu mujer,\\
aunque jamás se atrevieron ni a rozarla.\\
¿Por qué tanta violencia entonces, por qué la carnicería?\\
¿En el nombre de qué? ¿Del honor? ¿Por venganza?\\
Dime: ¿realmente lo pasaste tan mal en casa de Calipso?\\
¿Negarás que el amor de Circe era tan ardiente\\
como el mismísimo infierno que solo por ella recorriste?
\end{versopropio}

Y eso no es todo. Si de remordimientos se trata, Ulises tendrá que lamentar para siempre el acto que sigue a la carnicería de los pretendientes y que merece mención de honor en el libro universal de la infamia.
